\documentclass[12pt]{article}
\usepackage[utf8]{inputenc}
\usepackage[T1]{fontenc}
\usepackage{hyperref}
\title{Usarufa Language Resources for K-12 Curriculum Development}
\author{Compiled by Claude Sonnet 5 (claude-sonnet-5) for Andr\'as Kornai}
\date{}
\begin{document}
\maketitle
\section{Introduction}
Usarufa (ISO 639-3 usa) is a Kainantu-family language of the Eastern Highlands Province, Papua New Guinea, with about 1,200--1,300 speakers. Search found a grammar, New Testament portions and general reference pages, but no dictionary, primer or school materials.
\section{Primary Educational Resources}
None found in search results.
\section{Linguistic and Reference Materials}
\begin{itemize}
\item Bee, Darlene (1973), descriptive grammar of Usarufa, in McKaughan (ed.), \emph{The Languages of the Eastern Family of the East New Guinea Highland Stock}; also phonology papers by Bee and Glasgow (1962) and Bee (1964). Bibliographic record: \url{https://wals.info/refdb/record/Bee-1973}
\item Usarufa page of the New Guinea World project: \url{https://newguineaworld.linguistik.uzh.ch/families/trans-new-guinea/eastern-highlands-kratke-range/eastern-highlands/kainantu/west-kainantu/central-kainantu/usarufa}
\item ASJP wordlist for Usarufa: \url{https://asjp.clld.org/languages/USARUFA}
\item Wikipedia entry: \url{https://www.wikipedia.org/wiki/Usarufa}
\end{itemize}
\section{Religious and Community Resources}
\begin{itemize}
\item Usarufa Bible portions and New Testament on eBible.org (CC BY-NC-ND 4.0 according to the search summary): \url{https://ebible.org/details.php?id=usa}
\item Same text on png.bible: \url{https://www.png.bible/details.php?id=usa}
\item ANU repository catalogue record for a Usarufa New Testament: \url{https://openresearch-repository.anu.edu.au/items/f089bfef-d61b-403d-bd2c-91daacdc7202}
\end{itemize}
\section{Digital and Audio Resources}
None found in search results.
\section{Curriculum Development Context}
The language is not known to be taught in schools, and reportedly has no fluent speakers under about 25 (per Wikipedia, 2009). The scripture text and Bee's grammar are the only substantial written bases found; a school curriculum would need new materials.
\end{document}